\scanpage{71}
\uline{Digression} Épinglages de revêtements par \deleted{\unsure{cadres|cadre} fibrés \unsure{tangts|tangents}} rev. Kummériens des \deleted{\unsure{f|fi}}espaces tangents (cas relatif)
\par\noindent\rule{\linewidth}{0.4pt}\par

Soient
\[
 f:X\longrightarrow S \tag*{(1)}
\]
un morphisme lisse de dim relative 1, $t$ une section de $X$ sur $S$, $X^\ast=X\setminus t(S)$.

On considère des schémas \inserted{$X'$} finis sur $\deleted{S} X$, lisses sur $S$ de dim relative 1, munis d'une section $t'$ au dessus de $t$, et ``modérément ramifiés'' en $t'$ sur $X$, avec un degré de ramification $n$ -- ce qui signifie que
% slanted
\margin{\uline{NB} $n$ est premier aux car. résiduelles de $S$ !!!}
loc pour la top. étale sur $X,X'$, $X'$ est isom au voisinage de tt point $x'$ de $t'(S)$ à un $X$-schéma Kummérien de $X$ au voisinage de $x=p(x')$, de la forme
\margin{$p:X'\longrightarrow X$}
\[
 X'_0=\mathrm{Spec}\,\underline O_X[z]/(z^n-\varphi)
\]
où $\varphi$ est une équation locale du diviseur relatif $t(S)\subset X$ au voisinage de $x$.

Sous ces conditions, via l'hom norme locale, on définit un isom. canonique
\[
 T_{X'/\deleted{t}S,t'}^{\otimes n}\xrightarrow{\overset{k_{t'}}{\sim}}T_{X/S,t} \tag*{(2)}
\]
% archive sheet number: 70
